\chapter{Primary Decomposition}

\section*{PRIMARY IDEALS}

\begin{definition}
	An ideal $\mathfrak{q}\in A$ is \textcolor{orange}{primary} if $\mathfrak{q}\neq A$, and if $x\notin\mathfrak{q}\Rightarrow y^n\in\mathfrak{q}$ for some $n>0$. This also means $x\notin\sqrt{\mathfrak{q}}\Rightarrow y\in\mathfrak{q}$.
\end{definition}

\begin{proposition}{4.2}
	If $\sqrt{\mathfrak{a}}$ is maximal, then $\mathfrak{a}$ is primary.
	\begin{proof}
		Let $xy\in\mathfrak{a}$ but $x\notin\sqrt{\mathfrak{a}}$. We have $\sqrt{\mathfrak{a}}+(x)=A$, so $m+ax=1$ with $m^n\in\mathfrak{a}$. Therefore $m^n+bx=1$ and multiply $y$ to both sides.
	\end{proof}
\end{proposition}

\begin{proposition}{4.4}
	Let $\mathfrak{q}$ be a $\mathfrak{p}$-primary ideal, $x\in A$ an element. Then
	\begin{enumerate}[label=\textup{(\roman*)},ref=\thepropositiontemp(\roman*)]\crefalias{enumi}{proposition}
		\item if $x\in\mathfrak{q}$ then $(\mathfrak{q}:x)=(1)$;
		\item if $x\notin\mathfrak{q}$ then $(\mathfrak{q}:x)$ is $\mathfrak{p}$-primary;
		\item\label{4.4.3} if $x\notin\mathfrak{p}$ then $(\mathfrak{q}:x)=\mathfrak{q}$;
	\end{enumerate}
\end{proposition}

\begin{proposition}{4.8}
	Let $S$ be a multiplicatively closed subset of $A$, and let $\mathfrak{q}$ be a $\mathfrak{p}$-primary ideal.
	\begin{enumerate}[label=\textup{(\roman*)}]
		\item If $S\cap\mathfrak{p}\neq\varnothing$, then $S^{-1}\mathfrak{q}=S^{-1}A$.
		\item If $S\cap\mathfrak{p}=\varnothing$, then $S^{-1}\mathfrak{q}$ is $S^{-1}\mathfrak{p}$-primary and its contraction in $A$ is $\mathfrak{q}$.
	\end{enumerate}
\end{proposition}

\begin{definition}
	A \textcolor{orange}{primary decomposition} of an ideal $\mathfrak{a}$ in $A$ is an expression $\mathfrak{a}=\bigcap_{i=1}^n\mathfrak{q}_i$.
	
	If $\sqrt{\mathfrak{q}_i}$ are all distinct and $\mathfrak{q}_i\nsupseteq\bigcap_{j\neq i}\mathfrak{q}_j$, the primary decomposition is said to be \textcolor{orange}{irredundant}.
\end{definition}

\section*{UNIQUENESS}

\begin{proposition}{4.5}[1st uniqueness theorem]
	Let $\mathfrak{a}$ be a decomposable ideal and let $\mathfrak{a}=\bigcap_{i=1}^n\mathfrak{q}_i$ be a minimal primary decomposition of $\mathfrak{a}$. Let $\mathfrak{p}_i=\sqrt{\mathfrak{q}_i}$. Then $\mathfrak{p}_i$ are precisely the prime ideals which occur in the set of ideals $\sqrt{(\mathfrak{a}:x)}$, and hence are independent of the particular decomposition of $\mathfrak{a}$.
\end{proposition}

The prime ideals $\mathfrak{p}_i$ are said to belong to $\mathfrak{a}$, or to be \textcolor{orange}{associated} with $\mathfrak{a}$. The minimal elements of the set $\left\{\mathfrak{p}_1,\cdots,\mathfrak{p}_n\right\}$ are called the minimal or \textcolor{orange}{isolated} prime ideals belonging to $\mathfrak{a}$. The others are called \textcolor{orange}{embedded} prime ideals.

\begin{proposition}{4.9}
	Let $S$ be a multiplicatively closed subset of $A$ and let $\mathfrak{a}$ be a decomposable ideal. Let $\mathfrak{a}=\bigcap_{i=1}^n\mathfrak{q}_i$ be a minimal primary decomposition of $\mathfrak{a}$. Let $\mathfrak{p}_i=\sqrt{\mathfrak{q}_i}$ and suppose the $\mathfrak{q}_i$ numbered so that $S$ meets $\mathfrak{p}_{m+1},\cdots,\mathfrak{p}_n$ but not $\mathfrak{p}_1,\cdots\mathfrak{p}_m$. Then we have minimal primary decompositions
	\begin{equation*}
		S^{-1}\mathfrak{a}=\bigcap_{i=1}^mS^{-1}\mathfrak{q}_i,\qquad S(\mathfrak{a})=\bigcap_{i=1}^m\mathfrak{q}_i.
	\end{equation*}
\end{proposition}

A set $\Sigma$ of prime ideals belonging to $\mathfrak{a}$ is said to be \textcolor{orange}{isolated} if it satisfies the following condition: if $\mathfrak{p}^\prime$ is a prime ideal beloning to $\mathfrak{a}$ and $\mathfrak{p}^\prime\subseteq\mathfrak{p}$ for some $\mathfrak{p}\in\Sigma$, then $\mathfrak{p}^\prime\in\Sigma$. Thus, every isolated set contains an isolated prime ideal.

\begin{proposition}{4.10}[2nd uniqueness theorem]
	Let $\mathfrak{a}$ be a decomposable ideal, let $\mathfrak{a}=\bigcap_{i=1}^n\mathfrak{q}_i$ be a minimal primary decomposition of $\mathfrak{a}$, and let $\left\{\mathfrak{p}_{i_1},\cdots,\mathfrak{p}_{i_m}\right\}$ be an isolated set of prime ideals of $\mathfrak{a}$. Then the ideal $\mathfrak{q}_{i_1}\cap\cdots\cap\mathfrak{q}_{i_m}$ is independent of the decomposition.
\end{proposition}

\begin{corollary}{4.11}
	The isolated primary components are uniquely determined by $\mathfrak{a}$.
\end{corollary}

\section*{DEFINITIONS AND PROPOSITIONS FROM EXERCISES}

\begin{definition}
	Let $A$ be a ring and $\mathfrak{p}$ a prime ideal of $A$. The $n$th \textcolor{orange}{symbolic power} of $\mathfrak{p}$ is defined to be the ideal $\mathfrak{p}^{(n)}=S_\mathfrak{p}\left(\mathfrak{p}^n\right)$.
\end{definition}

\begin{definition}
	Let $M$ be an $A$-module. A submodule $P\subseteq M$ is called a \textcolor{orange}{prime submodule} if for any $a\in A$ and $x\in M$, $ax\in P$ implies $x\in P$ or $aM\subseteq P$.
	
	A submodule $Q\subseteq M$ is called a \textcolor{orange}{primary submodule} if $ax\in Q$ implies $x\in Q$ or $a^nM\subseteq Q$ for some $n>0$.
\end{definition}

\begin{definition}\label{primarymodule}
	Let $M$ be an $A$-module, $N$ a submodule of $M$. The \textcolor{orange}{ideal radial} $\operatorname{rad}_M(N)$ of $N$ in $M$ is defined to be $\operatorname{rad}_M(N)=\sqrt{(N:M)}$. The \textcolor{orange}{prime radical} $\operatorname{Rad}_M(N)$ of $N$ in $M$ is defined to be the intersection of all prime submodules containing $N$.
\end{definition}

\begin{proposition*}
	The ideal radical and the prime radical satisfy
	\begin{equation*}
		\left(\operatorname{Rad}_M(N):M\right)=\operatorname{rad}_M(N).
	\end{equation*}
	\begin{proof}
		$\left(\operatorname{Rad}_M(N):M\right)=\bigcap_{P\supseteq N}(P:M)$, so we need to proof $\bigcap_{P\supseteq N}(P:M)=\sqrt{(N:M)}$.
		
		If $N\subseteq P$, $(N:M)\subseteq(P:M)$. Since $(P:M)$ prime, $\sqrt{(N:M)}\subseteq(P:M)$ and therefore $\sqrt{(N:M)}\subseteq\bigcap_{P\supseteq N}(P:M)$.
		
		Let $a\in(P:M)$ for all $P\supseteq N$. For any $\mathfrak{p}\supseteq(N:M)$, let
		\begin{equation*}
			K_\mathfrak{p}(N)=\left\{m\in M:sm\in\mathfrak{p}M+N\text{ for some }s\in A-\mathfrak{p}\right\},
		\end{equation*}
		we have $N\subseteq K_\mathfrak{p}(N)$. If $(K_\mathfrak{p}(N):M)=\mathfrak{p}$, then $a\in\mathfrak{p}$ and therefore $a\in\sqrt{(N:M)}$.
		
		Let $x\in(K_\mathfrak{p}(N):M)$, then $xsm\in\mathfrak{p}M+N$ for all $m\in M$ and some $s\in A-\mathfrak{p}$. If $x\notin\mathfrak{p}$, $xs\in A-\mathfrak{p}$, and therefore $K_\mathfrak{p}(N)=M$, contradiction. Hence $(K_\mathfrak{p}(N):M)\subseteq\mathfrak{p}$.
		
		For any $x\in\mathfrak{p}$, $xM\subseteq\mathfrak{p}M\subseteq\mathfrak{p}M+N$. Therefore, $x\in(K_\mathfrak{p}(N):M)$ and so $\mathfrak{p}\subseteq(K_\mathfrak{p}(N):M)$.
	\end{proof}
\end{proposition*}

